
Practitioners selecting language models for safety-sensitive applications --- healthcare, legal information, public-facing dialogue --- routinely rely on trustworthiness benchmark scores: fairness audits, adversarial robustness rankings, alignment measurements. The implicit assumption is that a model ranked favorably on an in-distribution benchmark will remain favorable on out-of-distribution (OOD) conditions that differ in context or difficulty. This assumption is widely made but has not been empirically validated at scale across multiple trustworthiness dimensions.

Prior large-scale trustworthiness evaluations --- TrustLLM \citep{Sun2024TrustLLM}, DecodingTrust [Wang et al., 2023a], HELM \citep{Liang2022Holistic} --- report absolute performance scores on individual benchmarks. None asks whether in-distribution model rankings predict OOD model rankings across dimensions as a primary research question. The field's closest methodological precedent is Gevers and Daelemans [2026], who apply rank correlation with capability control to commonsense benchmarks (preprint; independently verified as plausible, but not confirmed via library access at submission time), not to trustworthiness dimensions with their distinctive latent-property versus adversarial-design structure.

A motivation for expecting \textit{differential} stability across dimensions comes from DecodingTrust [Wang et al., 2023a], which documents a two-model rank reversal: GPT-4 shows higher standard benchmark scores than GPT-3.5 but greater vulnerability to adversarial jailbreaking. This N = 2 observation suggested that adversarial benchmark design might disrupt rank preservation at scale, producing a meaningful gap between fairness and robustness stability. We test this prediction systematically.

This study fills the empirical gap with a focused analysis: partial Spearman $\rho$ (MMLU-controlled) between in-distribution and OOD model rankings across three trustworthiness benchmark pairs --- BBQ-Disambig $\rightarrow$ BBQ-Ambig (fairness), ANLI R1 $\rightarrow$ R3 (adversarial robustness), and OOD robustness --- for N = 16 LLMs from TrustLLM. The key methodological contribution is the capability control: raw Spearman $\rho$ is near-ceiling across all dimensions ($\rho \approx$ 0.978--0.984), masking dimension-specific structure. After partial correlation removing MMLU rank, the fairness-robustness differential becomes visible.

Three principal findings emerge. First, fairness rank ordering is highly stable across BBQ context splits after capability control (partial $\rho$ = 0.962, p < 0.0001, CI = [0.90, 1.00], N = 16). Second, contrary to the original hypothesis, adversarial robustness benchmark construction does not disrupt cross-split rank stability: both ANLI R1 $\rightarrow$ R3 ($\rho$ = 0.684, p = 0.007) and OOD robustness ($\rho$ = 0.868, p = 0.0001) show significantly positive partial $\rho$, with zero rank reversals across N = 13 models. The proposed adversarial disruption mechanism --- supported by the GPT-3.5/4 anecdote from DecodingTrust --- is falsified at scale. Third, the directional hypothesis that fairness stability exceeds robustness stability receives partial support ($\Delta\rho$ = 0.192, Fisher z p = 0.024, N = 13), though the preregistered criterion ($\Delta\rho \geq$ 0.200) was not met, falling short by 0.008.

This work makes three contributions:

\begin{enumerate}
\item \textbf{Empirical baseline.} First computation of partial Spearman $\rho$ (MMLU-controlled) between in-distribution and OOD trustworthiness benchmark model rankings across 16 LLMs, establishing that fairness ($\rho$ = 0.962) and adversarial robustness ($\rho$ = 0.684--0.868) rankings are each highly stable across distribution shifts.
\end{enumerate}

\begin{enumerate}
\item \textbf{Methodological demonstration.} MMLU-controlled partial Spearman $\rho$ is a tractable, data-efficient operationalization of trustworthiness benchmark predictive validity using published scores only --- no new data collection required.
\end{enumerate}

\begin{enumerate}
\item \textbf{Theoretical revision.} The adversarial disruption hypothesis is falsified. Both fairness and robustness are stable latent model properties. The directional fairness advantage ($\Delta\rho$ = 0.192, p = 0.024) lacks a confirmed mechanism, setting an open research agenda.
\end{enumerate}

